\chapter{Introduction}
\label{cap:introduction}

This report documents the design, implementation and validation of the
\emph{Talent Intelligence and Language Verification Platform}, a web
application built during the BlendEd program in partnership with adidas
Global Business Services (GBS) Porto, in the context of the LEI-PROJ
(PESTI) project unit of the Bachelor's Degree in Computer Engineering at
the Instituto Superior de Engenharia do Porto.

The platform addresses an everyday tension faced by recruitment teams that
operate under European data protection law. On one hand, hiring volumes are
high and recruiters spend a substantial share of their time triaging
\textit{Curricula Vitae}, screening for language proficiency and tracking
status changes across spreadsheets and email threads. On the other hand,
the General Data Protection Regulation (\gls{gdpr}) constrains how long
candidate data may be retained after a process closes, which forces
periodic deletions and turns institutional memory into a liability rather
than an asset. The work presented here proposes a structured talent pool
backed by automated CV parsing, job-anchored matching, and a Large Language
Model (\gls{llm}) based assessment module that together help HR teams
extract value from candidate data within the time window the law allows,
and re-engage promising candidates in a compliant way.

%----------------------------------------------------------------------------------------

\section{Context}
\label{sec:context}

The BlendEd program is a partnership between Portuguese higher education
institutions and adidas GBS Porto in which student teams design and
implement software solutions for problems sourced from the company's day-to-day
operations. The team behind this project (two members) was assigned to the
HR / Talent Acquisition area at adidas GBS Porto with a deliberately broad
brief: to identify recurring frictions in the recruitment workflow and to
propose a tool that improves them.

Discovery sessions with the client surfaced a pattern that runs across most
of their processes. Candidate data is rich and heterogeneous (CVs in
multiple formats, scattered emails, spreadsheets per campaign, ad-hoc
language self-assessments) but it is not stored in a way that supports
analysis or reuse. The same information has to be re-extracted from CVs
every time a new opening is published, recruiters cannot easily tell
whether a candidate has applied before, and there is no consistent way to
verify language proficiency before a first interview. The six-month
retention rule imposed by \gls{gdpr} amplifies all of these problems: the
moment a process closes, the institutional memory it generated starts a
countdown.

This project is therefore positioned at the intersection of three concerns
that recruitment software rarely addresses simultaneously: HR productivity
(reducing manual CV analysis time and giving recruiters analytics-ready
views), candidate experience (proactive communication and feedback paths
for candidates not selected) and regulatory compliance (a deletion-aware
data model and a contact-history audit trail).

%----------------------------------------------------------------------------------------

\section{Problem Statement}
\label{sec:problem}

The HR teams interviewed during the discovery phase reported nine concrete
pain points, summarised in
Table~\ref{tab:pain-points-overview}. These were validated through repeated
client meetings and an instrumented survey with one question per pain
point.

\begin{table}[h]
\centering
\caption[Pain points identified during client discovery]{Pain points
identified during client discovery, as documented in
\textit{TalentPool-client\_pain\_points.md} (29\textsuperscript{th} April 2026).}
\label{tab:pain-points-overview}
\begin{tabular}{p{0.05\linewidth}p{0.85\linewidth}}
\toprule
\textbf{\#} & \textbf{Pain point} \\
\midrule
1 & No structured way to organise candidate information; spreadsheets are still common \\
2 & Difficult to know whether a candidate has applied before, made worse by the GDPR six-month rule \\
3 & No first-pass screening for language or other skills before the first interview \\
4 & No structured channel to contact candidates (status changes, opportunities, campaigns) \\
5 & Significant time spent extracting basic information from each CV manually \\
6 & No quick, organised view of the recruitment pipeline \\
7 & No data-driven way to decide where to focus future recruitment efforts \\
8 & Existing tools do not feel like they reduce repetitive work \\
9 & No structured process to keep promising-but-rejected candidates engaged \\
\bottomrule
\end{tabular}
\end{table}

The problem this project addresses can therefore be stated as follows:

\begin{quote}
\textit{How can a small HR team manage a structured pool of candidates --- including
CV-derived facts, language proficiency, application history and outbound
communications --- in a way that is efficient under high volume, compliant
with GDPR retention obligations, and useful both for individual hiring
decisions and for strategic planning?}
\end{quote}

%----------------------------------------------------------------------------------------

\section{Objectives}
\label{sec:objectives}

The project pursues seven objectives, derived directly from the pain points
above and refined during sprint planning with the client:

\begin{enumerate}
	\item \textbf{Centralise candidate data.} Replace ad-hoc spreadsheets
	with a relational, GDPR-aware talent pool storing CV-derived facts,
	application history, languages, skills, notes and assessment results.

	\item \textbf{Automate CV ingestion.} Parse uploaded CVs with a
	\gls{llm} pipeline (Groq Llama~3.3 70B as primary, OpenAI GPT-4o as
	fallback) that extracts structured experience, education, languages,
	skills and a per-experience field-of-work tag.

	\item \textbf{Provide job-anchored matching.} Score candidates against
	specific job descriptions using a rule-based fit algorithm
	(\textit{computeJobFit}) that combines field of work, seniority,
	required and preferred skills, languages and education, and that
	exposes per-criterion explanations.

	\item \textbf{Verify language proficiency.} Offer an AI-driven
	conversational interview that returns a Common European Framework of
	Reference (\gls{cefr}) level along with grammar, vocabulary and
	fluency sub-scores, plus a separate technical-skill interview mode.

	\item \textbf{Empower HR communication.} Build a notification and
	campaign system with full interaction history per candidate, including
	manually composed emails, status changes and promotional outreach.

	\item \textbf{Make the pipeline measurable.} Provide both built-in
	dashboards and a custom widget builder so HR can observe pipeline
	health and slice the data along whitelisted dimensions.

	\item \textbf{Respect compliance and security.} Enforce an
	authentication-and-authorisation model that gates HR-only operations
	server-side, and adopt a deletion-aware data model that aligns with
	the six-month GDPR retention rule.
\end{enumerate}

%----------------------------------------------------------------------------------------

\section{Approach}
\label{sec:approach}

The project followed an incremental, Scrum-based delivery cycle over the
academic year, with the author acting as Product Owner and the second team
member, Stratos Demertzoglou, owning the AI Interviewer sidecar (a Python
\textit{FastAPI} service) end to end. All other modules --- frontend,
Next.js API routes, database schema, CV parsing pipeline, matching engine,
notifications, analytics and CI/CD --- were designed and implemented by
the author.

The technological choices were guided by the brief: a small two-person
team, a tight academic deadline, and the need for production-grade
deployment. The platform is implemented as a Next.js 16 application
written in TypeScript, deployed on Vercel, and backed by a Supabase
PostgreSQL database with Supabase Auth (Google OAuth). The codebase
follows an onion / clean architecture (\gls{onion}) in which the Domain
layer has zero external dependencies and Use Cases never import from
infrastructure --- a discipline that proved decisive when the AI
Interviewer module had to be slotted in as a separate process. Validation
combines automated unit and integration testing with Vitest, manual
acceptance testing with the client, and a structured feedback survey
mapped one-to-one onto the nine pain points.

The contributions claimed by this work are:

\begin{itemize}
	\item A reference architecture for an AI-assisted recruitment platform
	that respects GDPR retention rules without sacrificing analytical
	value.
	\item A reproducible job-anchored fit algorithm that is auditable
	per-criterion (as opposed to opaque embedding-based ranking),
	implemented as a pure function with full unit-test coverage.
	\item A dual-mode AI interviewer (technical skill validation and
	CEFR-grade language assessment) integrated into a Next.js application
	through a hardened HMAC-signed token contract.
	\item A documented project-management experience report from a
	two-person team using Scrum, with the author serving as Product Owner.
\end{itemize}

%----------------------------------------------------------------------------------------

\section{Team Composition and Personal Contribution}
\label{sec:team}

The project was developed by a team of two:

\begin{description}
	\item[Fernando Regalado Lobo Ribeiro (author):] Product Owner;
	full-stack engineer; responsible for client discovery, requirements,
	architecture, database design, the entire Next.js application
	(frontend, API routes, use cases, repositories, AI parsing pipeline,
	matching engine, notifications, analytics, security, CI/CD,
	deployment), automated testing and the writing of this report.

	\item[Stratos Demertzoglou:] AI engineer; responsible for the AI
	Interviewer FastAPI service: the conversational orchestration loop,
	the CEFR evaluator, the system prompts for both technical and
	language modes, and the integration contract consumed by the
	Next.js side.
\end{description}

The split was negotiated explicitly during sprint zero, codified in the
product backlog as a module boundary, and re-confirmed at every sprint
review. Because the author was simultaneously Product Owner and the main
implementer, the role had to be exercised carefully: backlog ordering had
to remain client-driven and not author-driven, and integration risk
between the two modules was mitigated through a documented HTTP contract
and a shared session-token scheme described in
Chapter~\ref{cap:design}.

%----------------------------------------------------------------------------------------

\section{Document Structure}
\label{sec:structure}

The remainder of this report is organised as follows.
Chapter~\ref{cap:state-of-the-art} reviews the state of the art in
recruitment software, AI-assisted CV parsing, conversational language
assessment and the regulatory landscape that shapes them.
Chapter~\ref{cap:analysis} analyses the problem in depth, formalises the
requirements and describes the project-management approach and the team
dynamics.
Chapter~\ref{cap:design} presents the solution design: the architecture,
domain model, database schema, matching algorithm and the AI Interviewer
contract.
Chapter~\ref{cap:implementation} walks through the implementation,
covering the technology stack, the most relevant features and the
trade-offs encountered along the way.
Chapter~\ref{cap:validation} reports verification and validation results,
including automated test coverage, security review and the client-facing
survey.
Chapter~\ref{cap:conclusions} closes the report with conclusions, an
honest assessment of limitations and a roadmap for future work.
The bibliography and a set of appendices --- detailed database schema,
full API reference, and the survey instrument --- complete the document.
